\subsection{Universal covering of a deloopable topological group}

Let G be a locally path-connected topological group. The translations of G are homeomorphisms (TG, III, p.~2). For the space G to be deloopable (IV, p.~340, def. 2), it is necessary and sufficient that G possess the following property:

There exists a neighborhood V of the identity element e of G such that the image of the homomorphism from \(\pi_{1}(V,e)\) into \(\pi_{1}(G,e)\) deduced from the canonical injection is reduced to the identity element.

PROPOSITION 6. --- Let G be a connected and deloopable topological group. There exists a topological group \(\widetilde{\mathbf{G}}\) with identity element \(\widetilde{e}\) , and a continuous homomorphism \(p\colon \widetilde{\mathbf{G}}\to \mathbf{G}\) satisfying the following assertions:

\begin{enumerate}
\def\labelenumi{\alph{enumi})}
\item
  The space \(\widetilde{G}\) is simply connected and simply path-connected. Equipped with the map p, the pointed space ( \(\widetilde{G},\widetilde{e}\) ) is a universal covering of the pointed space (G,e).
\item
  The kernel N of p is a discrete subgroup of \(\widetilde{G}\), contained in the center of \(\widetilde{G}\) . The homomorphism \(\mathrm{N}\to\mathrm{Aut}_{\mathrm{G}}(\widetilde{\mathrm{G}})\) which sends \(n\in N\) to the right translation in \(\widetilde{G}\) , is a group isomorphism. The homomorphism from \(\pi_{1}(\mathrm{G})\) into N which sends \(\gamma\in\pi_{1}(\mathrm{G})\) to the unique element n of N such that \(\widetilde{e}\cdot\gamma=n\) , is an isomorphism of groups.
\item
  If G' is a topological group, with identity element e', and if p': G' → G is a homomorphism of groups which makes G' a covering of G, the unique continuous map u: \(\widetilde{G} \rightarrow G'\) such that \(u(\widetilde{e}) = e'\) and \(p' \circ u = p\) is a homomorphism of groups. Equipped with the map u, \((\widetilde{G}, \widetilde{e})\) is a universal covering of \((G', e')\) .
\end{enumerate}

By IV, p.~342, theorem 1, there exists a covering \((\widetilde{\mathbf{G}},p)\) of G, simply connected and simply path-connected, Galois with group \(\pi_1(\mathbf{G})^\circ\) , and a point \(\widetilde{e}\) of \(p^{-1}(e)\) such that the pointed covering \((\widetilde{\mathbf{G}},\widetilde{e})\) is a universal covering of \((\mathbf{G},e)\) .

By IV, p.~375, prop. 4 and IV, p.~377, prop. 5, there exists on the space \(\widetilde{\mathrm{G}}\) a unique group structure compatible with its topology for which \(p\) is a homomorphism of groups and for which \(\widetilde{e}\) is an identity element. It follows from prop. 4 that the group \(\widetilde{\mathrm{G}}\) satisfies the assertions \(a)\) and \(b)\) .

Let us prove assertion \(c\). Let \(\mathbf{G}'\) be a topological group and let \(p'\colon \mathbf{G}' \to \mathbf{G}\) be a homomorphism of groups which makes \(\mathbf{G}'\) a covering of \(\mathbf{G}\) . Let \(e'\) be the identity element of \(\mathbf{G}'\) . The existence and uniqueness of a map \(u\colon \widetilde{\mathbf{G}} \to \mathbf{G}'\) such that \(p = p' \circ u\) and \(u(\widetilde{e}) = e'\) follow from the fact that \((\widetilde{\mathbf{G}}, \widetilde{e})\) is a universal covering of the pointed space \((\mathbf{G}, e)\) . Since the maps \(p\) and \(p'\) are surjective homomorphisms of groups, \(u\) is a homomorphism of groups. Equipped with the map \(u\) , \(\widetilde{\mathbf{G}}\) is a covering of \(\mathbf{G}'\) (I, p.~81, prop. 7), hence \((\widetilde{\mathbf{G}}, \widetilde{e})\) is a universal covering of \((\mathbf{G}', e')\) (IV, p.~345, cor. 2 of th. 1).

Under the hypotheses of the proposition, one will say, by abuse of language, that \(\tilde{G}\) is a universal covering of \(G\) .

Example. --- Proposition 6 applies in particular when G is a connected Lie group over \(\mathbf{R}\) or \(\mathbf{C}\). There then exists on \(\widetilde{\mathbf{G}}\) a unique analytic manifold structure such that the projection \(p\colon\widetilde{\mathbf{G}}\to\mathbf{G}\) is an étale morphism of analytic manifolds (VAR R, §1, 5.8.2, p.~48). For this manifold structure, \(\widetilde{\mathbf{G}}\) is a Lie group (LIE, III, p.~113, corollary).

Scholium. --- Let G be a connected and deloopable topological group, let e be its identity element. Let \(\widetilde{G}\) be a simply path-connected topological group, with identity element \(\widetilde{e}\) , and let \(p\colon\widetilde{G}\to G\) be a homomorphism of groups which makes \(\widetilde{G}\) a universal cover of G. Let N denote the kernel of p.

Right translation \(\delta_{h}\) (resp. left translation) by an element \(h \in N\) is a G-automorphism of the principal covering \(\widetilde{G}\) and the mapping \(h \mapsto \delta_{h}\) is an isomorphism from the group N onto the group of automorphisms of this principal covering.

For every subgroup K of N, the map \(p'\colon\widetilde{G}/K\to G\) deduced from p is a Galois covering of G, and \(\operatorname{Aut}_{\mathrm{G}}(\widetilde{\mathrm{G}}/\mathrm{K})\) is identified with the group N/K. When one identifies the groups N and \(\pi_{1}(\mathrm{G})\) , the group \(p'_{*}(\pi_{1}(\widetilde{\mathrm{G}}/\mathrm{K}))\) is identified with the subgroup K of N. Moreover, \(\widetilde{G}\) is a covering of \(\widetilde{G}/K\) , since G is locally connected (I, p.~81, prop. 7).

Conversely, every nonempty connected covering E of G is G-isomorphic to a covering of this type (I, p.~113, th. 3 and I, p.~111, prop. 10). Indeed, consider a point x of the fiber \(E_{e}\) and let f be the unique homomorphism from \(\widetilde{G}\) into E which maps \(\widetilde{e}\) to x. Endowed with f, \(\widetilde{G}\) is a Galois covering of E; the subgroup \(\mathrm{Aut}_{\mathrm{E}}(\widetilde{\mathrm{G}})\)

of \(\mathrm{Aut}_{\mathrm{G}}(\widetilde{\mathrm{G}})\) is identified with \(f^{-1}(x)\) which is therefore a subgroup of N. Consequently, \(f\) induces a G-isomorphism of \(\widetilde{\mathrm{G}} / f^{-1}(x)\) onto E. By transport of structure, this yields a topological group structure on E for which \(x\) is the identity element and for which the map \(f\) is a surjective homomorphism. The projection of the G-space E is then a group homomorphism, and the composition law of E is therefore the composition law defined by Prop. 4 of IV, p.~375.

Let \((\mathrm{E},q)\) and \((\mathrm{E}^{\prime},q^{\prime})\) be connected coverings of G, let x be a point of \(E_{e}\) and \(x^{\prime}\) be a point of \(E_{e}^{\prime}\). For there to exist a G-morphism from E into \(E^{\prime}\) , it is necessary and sufficient that \(p_{*}(\pi_{1}(\mathrm{E},x))\) be contained in \(p_{*}^{\prime}(\pi_{1}(\mathrm{E}^{\prime},x^{\prime}))\) . If this condition is satisfied, there then exists a unique G-morphism \(f\colon E\to E^{\prime}\) such that \(f(x)=x^{\prime}\) (III, p.~311, cor. 2 of prop. 1). If one equips E and \(E^{\prime}\) with group composition laws for which q and \(q^{\prime}\) are homomorphisms and x and \(x^{\prime}\) neutral elements, the map f is a group homomorphism. Indeed, f identifies with the canonical homomorphism. \(\widetilde{\mathrm{G}}/p_{*}(\pi_{1}(\mathrm{E},x))\to\widetilde{\mathrm{G}}/p_{*}^{\prime}(\pi_{1}(\mathrm{E}^{\prime},x^{\prime}))\) .

PROPOSITION 7. --- For two connected and deloopable topological groups to be locally isomorphic, it is necessary and sufficient that their universal covering groups be isomorphic as topological groups.

A connected and deloopable topological group is locally isomorphic to its universal covering (IV, p.~369, prop. 1); the condition is therefore sufficient. It is necessary by corollary 2 of proposition 2 (IV, p.~372), since the universal covering of a connected and deloopable topological group is simply connected (IV, p.~379, prop. 6).

PROPOSITION 8. --- Let G be a connected and deloopable topological group and let \(\widetilde{\mathbf{G}}\) be a universal covering of G. Let V be a connected open neighborhood of the identity element e of G such that the image of the canonical homomorphism from \(\pi_1(\mathrm{V} \cdot \mathrm{V}, e)\) into \(\pi_1(\mathrm{G}, e)\) is reduced to the identity element. Let F(V, r) denote the group defined by the generating set V and the set r of relators \(xyz^{-1}\), where \((x, y, z)\) ranges over the set of elements of V × V × V such that \(xy = z\) . Denote by j: V → F(V, r) the canonical map.

There exists a unique isomorphism \(f\) of the group \(\mathrm{F}(\mathrm{V},\mathbf{r})\) onto \(\widetilde{\mathrm{G}}\) such that \(f \circ j\) is a lift to \(\widetilde{\mathrm{G}}\) of the canonical injection from V into G.

The set V · V is connected and open, hence locally path-connected. Let p be the projection of \(\widetilde{G}\). There exists a continuous section s of p over V · V such that \(s(e) = \widetilde{e}\) , where \(\widetilde{e}\) denotes the identity element of \(\widetilde{G}\) (III, p.~308, prop. 1). Set \(\widetilde{V} = s(V)\) ; the set \(\widetilde{V}\) is a connected open neighborhood of \(\widetilde{e}\) in \(\widetilde{G}\) and s is a homeomorphism of V · V onto \(\widetilde{V} \cdot \widetilde{V}\) . The maps \((x, y) \mapsto s(x)s(y)\) and \((x, y) \mapsto s(xy)\) are liftings to \(\widetilde{G}\) of the map \((x, y) \mapsto xy\) from V × V into G that coincide at \((e, e)\) ; since V × V is connected, they coincide on V × V (I, p.~34, cor. 1). Moreover, if \((\widetilde{x}, \widetilde{y}, \widetilde{z}) \in \widetilde{V} \times \widetilde{V} \times \widetilde{V}\) , the conditions \(\widetilde{x}\widetilde{y} = \widetilde{z}\) and \(p(\widetilde{x})p(\widetilde{y}) = p(\widetilde{z})\) are equivalent. The existence of an isomorphism \(f: F(V, r) \to \widetilde{G}\) then follows from Corollary 3 (IV, p.~372) of Proposition 2.

Let \(g\) be an isomorphism of \(\mathrm{F}(\mathrm{V},\mathbf{r})\) satisfying the conditions of the proposition. The equality \(e \cdot e = e\) entails that \(eee^{-1} \in \mathbf{r}\) , so that \(j(e)\) is the identity element of \(\mathrm{F}(\mathrm{V},\mathbf{r})\) . Since V is connected, \(g \circ j\) is the unique continuous lift to \(\widetilde{\mathrm{G}}\) of the canonical injection of V into G. Consequently, \(f\) and \(g\) coincide on \(j(\mathrm{V})\) . Since \(j(\mathrm{V})\) generates the group \(\mathrm{F}(\mathrm{V},\mathbf{r})\) , \(f = g\) .

